\documentclass[10pt,a4paper]{amsart}
\usepackage[margin=19mm]{geometry}
\usepackage{amsmath,amssymb,mathtools}
\usepackage[scr=rsfs,cal=euler]{mathalfa}
\usepackage{xfrac}
\usepackage[unicode,hidelinks,bookmarks=false]{hyperref}
\newcommand{\ie}{i.e.\ }
\newcommand{\cf}{cf.\ }
\newcommand{\resp}{respectively\ }
\newcommand{\Subsection}{\subsection}
\providecommand{\nobreakdash}{\nobreak\hbox{-}}
\setlength{\emergencystretch}{2em}
\multlinegap0pt
\usepackage{bm}
\usepackage{amssymb}
\usepackage{mathtools}
\usepackage{xcolor}
\usepackage[shortlabels]{enumitem}
\newtheorem{theorem}{Theorem}
\newcommand{\HU}{\mathrm{HU}}
\newcommand{\Sh}{\operatorname{Sh}}
\newcommand{\Var}{\operatorname{Var}}
\newcommand{\Vol}{\operatorname{Vol}}
\newcommand{\cM}{\mathcal M}
\title{Higher-order hyperuniformity of random measures}
\author{Michael Bj\"orklund}
\date{Source: arXiv:2609.07502v1 (CC BY 4.0)}
\begin{document}
\maketitle

\noindent\emph{Source and licence.} Higher-order hyperuniformity of random measures. Version arXiv:2609.07502v1 (CC BY 4.0), distributed by arXiv under the Creative Commons Attribution 4.0 International licence, http://creativecommons.org/licenses/by/4.0/, as recorded in the arXiv licence metadata for this exact version and shown on the abstract page https://arxiv.org/abs/2609.07502v1. Full licence notice: \texttt{licenses/arxiv-2609.07502-CC-BY-4.0.txt}.\par\smallskip

\noindent\textit{Editorial scope.} Counted unit: the article's theorem thm:cap-codim-one-separation (source0.tex lines 4122--4132). The article's complete original proof of it is delivered with it (source0.tex lines 4134--4531, one \texttt{proof} environment of 398 lines). No mathematics is written, repaired or re-derived: the statement and the proof are the article's own lines, in the article's own notation, and no retained line is substituted or re-flowed.  Retained uncounted as the source-local context the proof needs: lines 1121--1125; lines 3869--3874; lines 4086--4094.  Nothing was omitted from any retained span, so the package carries no omission record.  No retained line contains a citation, so the wrapper carries no bibliography and no bibliography entry.  No retained line contains a figure environment, an image-inclusion command or drawing code, so no image is delivered and the package declares no asset dependency.  Editorial formatting only, in addition: an AMS \texttt{amsart} preamble that restates the article's own declarations and macros used by the retained lines, namely the environments \texttt{theorem} on separate counters, together with the standard packages those lines need.  The retained environments print in this document as Theorem 1 in order of appearance, whereas the article prints the same items with the numbering of its own full document.  The complete native source of the article as distributed in the arXiv e-print for this exact version is bundled unchanged as \texttt{sources/209634/source0.tex}.
\begin{theorem}[Geometric--spectral equivalence]\label{thm:geometric-spectral}
Let $\mu$ be a $V$-invariant probability measure on $\cM(V)$ which is
locally $L^{2k}$-integrable.  Then $\mu$ is geometrically $k$-hyperuniform
if and only if it is spectrally $k$-hyperuniform.
\end{theorem}

\begin{equation}\label{eq:cap-diffraction-formula}
  \sigma_{W'}
  =\frac{1}{\operatorname{covol}(\Gamma)^2}
   \sum_{(\xi,\eta)\in\Gamma^\perp\setminus\{0\}}
     |\widehat{\chi}_{W'}(\eta)|^2\,\delta_\xi.
\end{equation}

\begin{equation}\label{eq:cap-isolated-derived-window-statistic}
  T_k[f,\varphi]
  =c_{\mathbf m}\,
    S_{W[\boldsymbol\gamma]}
      (\tau_{b_{\mathbf m,\boldsymbol\gamma}}f),
  \qquad
  b_{\mathbf m,\boldsymbol\gamma}
  :=\frac1k\sum_{r=1}^j m_r\gamma_{r,V}.
\end{equation}

\begin{theorem}[Codimension-one separation]
\label{thm:cap-codim-one-separation}
There exists a Euclidean cut-and-project scheme
\[
  (\mathbb R^3,\mathbb R,\Gamma)
\]
and an interval window $W\subset\mathbb R$ such that
\begin{equation}\label{eq:cap-codim-one-separation}
  \mu_W\in\HU_1\setminus\HU_2.
\end{equation}
\end{theorem}

\begin{proof}
\smallskip\noindent\emph{The lattice and the window cancellation.}
Choose integers $M_n\geq4$ such that
\begin{equation}\label{eq:cap-codim-one-growth}
  M_n\mid M_{n+1},
  \qquad
  M_{n+1}/M_n\in M_n\mathbb Z,
  \qquad
  M_{n+1}\geq M_n^{12}.
\end{equation}
For example, one may take $M_1=4$ and $M_{n+1}=M_n^{12}$.  Set
\begin{equation}\label{eq:cap-codim-one-alpha}
  \alpha_j:=\sum_{n\geq1}M_n^{-j},
  \qquad j=1,2,3,
  \qquad
  \alpha:=(\alpha_1,\alpha_2,\alpha_3)^T.
\end{equation}
Then $0<\alpha_1<1$.  Each $\alpha_j$ is irrational.  Indeed, with
\begin{equation}\label{eq:cap-codim-one-prefixes-intro}
  \alpha_{j,n}:=\sum_{r\leq n}M_r^{-j}
  =\frac{P_{j,n}}{M_n^j},
\end{equation}
one has
\(
  0<\alpha_j-\alpha_{j,n}\leq2M_{n+1}^{-j}=o(M_n^{-j}).
\)
If $\alpha_j=a/b$ were rational, then the distinct rationals $a/b$ and
$P_{j,n}/M_n^j$ would differ by at least $(bM_n^j)^{-1}$, a contradiction for
large $n$.

Let $e_3=(0,0,1)^T$ and put
\begin{equation}\label{eq:cap-codim-one-G}
  G:=
  \begin{pmatrix}
    -I_3+e_3\alpha^T & e_3\\
    \alpha^T & 1
  \end{pmatrix}.
\end{equation}
A direct multiplication gives
\begin{equation}\label{eq:cap-codim-one-Gdual}
  G^{-T}
  =
  \begin{pmatrix}
    -I_3 & \alpha\\
    e_3^T & 1-\alpha_3
  \end{pmatrix}.
\end{equation}
The displayed matrix has determinant $-1$, hence $\det G=-1$ and
$\operatorname{covol}(\Gamma)=1$.  Define
\[
  \Gamma:=G\mathbb Z^4<\mathbb R^3\times\mathbb R,
  \qquad
  \Lambda:=\Gamma^\perp=G^{-T}\mathbb Z^4.
\]
For $m\in\mathbb Z^3$ and $r\in\mathbb Z$,
\[
  G\binom mr
  =\bigl(-m+e_3(\alpha\cdot m+r),\,\alpha\cdot m+r\bigr).
\]
If its physical component vanishes, then $m_1=m_2=0$ and
$r-(1-\alpha_3)m_3=0$.  Since $\alpha_3$ is irrational, $m_3=r=0$.
Thus the physical projection of $\Gamma$ is injective.  Its internal
projection contains $\mathbb Z+\alpha_3\mathbb Z$, hence is dense.  Therefore
$(\mathbb R^3,\mathbb R,\Gamma)$ is a Euclidean cut-and-project scheme.

Writing an element of $\Lambda$ as $(\xi,\eta)$ and the corresponding integer
vector as $(p_1,p_2,p_3,q)$, Equation~\eqref{eq:cap-codim-one-Gdual} gives
\begin{equation}\label{eq:cap-codim-one-dual-coordinates}
  \xi_j=q\alpha_j-p_j,
  \qquad
  \eta=p_3+(1-\alpha_3)q=q-\xi_3.
\end{equation}
The physical projection of $\Lambda$ is injective: if $\xi=0$, then
$q\alpha_1=p_1$, hence $q=0$ and then $p=0$.

Use the lattice vector
\begin{equation}\label{eq:cap-codim-one-primal-vectors}
  \gamma_1=(-e_1+\alpha_1e_3,\alpha_1).
\end{equation}
The remaining dual coordinate satisfies $\eta+\xi_3\in\mathbb Z$, which
produces the Fourier cancellation for the original interval.  Take
\begin{equation}\label{eq:cap-codim-one-window}
  W=[-1/2,1/2].
\end{equation}
Since $\eta+\xi_3=q\in\mathbb Z$, every dual point with $\eta\neq0$
satisfies
\begin{equation}\label{eq:cap-codim-one-window-cancellation}
  |\widehat{\chi}_W(\eta)|
  =\frac{|\sin(\pi\xi_3)|}{\pi|\eta|}
  \leq\frac{|\xi_3|}{|\eta|}.
\end{equation}
If $\eta=0$, irrationality of $\alpha_3$ forces $q=p_3=0$, and every
nonzero such dual point has $\xi\in\mathbb Z^3\setminus\{0\}$.  Hence these
points never occur in $B_\varepsilon^*$ once $\varepsilon<1$.

\smallskip\noindent\emph{Separation estimate for small physical frequencies.}
Put
\begin{equation}\label{eq:cap-codim-one-tails}
  t_{j,n}:=\alpha_j-\alpha_{j,n}.
\end{equation}
The growth assumptions imply
\begin{equation}\label{eq:cap-codim-one-tail-bounds}
  M_{n+1}^{-j}\leq t_{j,n}\leq2M_{n+1}^{-j}.
\end{equation}
Moreover $P_{j,n}$ is coprime to $M_n$.  Indeed, $P_{j,1}=1$, and if
$R_n:=M_{n+1}/M_n$, then
\begin{equation}\label{eq:cap-codim-one-prefix-recurrence}
  P_{j,n+1}=R_n^jP_{j,n}+1.
\end{equation}
Every prime divisor of $M_{n+1}=M_nR_n$ divides $R_n$, because
$M_n\mid R_n$.  Hence $P_{j,n+1}\equiv1$ modulo every prime divisor of
$M_{n+1}$.

Set
\begin{equation}\label{eq:cap-codim-one-qA}
  q_n:=M_n^3,
  \qquad
  A_n:=\frac{M_{n+1}}{8M_n}.
\end{equation}
Write $\operatorname{dist}_\infty(\cdot,\mathbb Z^3)$ for distance in the
sup norm.  If $0<|q|\leq A_n$ and $q$ is not divisible by $q_n=M_n^3$, choose
$r\in\{0,1,2\}$ such that
\[
  M_n^r\mid q,
  \qquad
  M_n^{r+1}\nmid q.
\]
Write $q=M_n^r u$.  Since
$\alpha_{r+1,n}=P_{r+1,n}/M_n^{r+1}$ and
$\gcd(P_{r+1,n},M_n)=1$,
\[
  q\alpha_{r+1,n}=\frac{uP_{r+1,n}}{M_n}.
\]
Because $M_n\nmid u$, this is not an integer.  After cancellation its
denominator is at most $M_n$, so its distance from $\mathbb Z$ is at least
$1/M_n$.  Therefore
\begin{equation}\label{eq:cap-codim-one-prefix-separation}
  \operatorname{dist}_\infty
  \bigl(q(\alpha_{1,n},\alpha_{2,n},\alpha_{3,n}),\mathbb Z^3\bigr)
  \geq\frac1{M_n}.
\end{equation}
For $|q|\leq A_n$, the tail bound gives
\begin{equation}\label{eq:cap-codim-one-tail-perturbation}
  \|q(t_{1,n},t_{2,n},t_{3,n})\|_\infty
  =|q|t_{1,n}
  \leq\frac1{4M_n}.
\end{equation}
For $q\in\mathbb Z$ with
$\operatorname{dist}_\infty(q\alpha,\mathbb Z^3)<1/2$, let $p_j(q)$ be the
unique nearest integer to $q\alpha_j$ and put
\begin{equation}\label{eq:cap-codim-one-return-coordinate-definition}
  \xi_j(q):=q\alpha_j-p_j(q),\qquad j=1,2,3.
\end{equation}
If the hypotheses on the left-hand side of
\eqref{eq:cap-codim-one-return-isolation} hold and $q\notin q_n\mathbb Z$,
then \eqref{eq:cap-codim-one-prefix-separation} and
\eqref{eq:cap-codim-one-tail-perturbation} give
\[
  \operatorname{dist}_\infty(q\alpha,\mathbb Z^3)
  \geq \frac1{M_n}-\frac1{4M_n}
  =\frac3{4M_n},
\]
a contradiction.  Hence
\begin{equation}\label{eq:cap-codim-one-return-isolation}
  0<|q|\leq A_n,
  \qquad
  \operatorname{dist}_\infty(q\alpha,\mathbb Z^3)<\frac1{2M_n}
  \quad\Longrightarrow\quad
  q\in q_n\mathbb Z.
\end{equation}
If $q=kq_n$ and $|q|\leq A_n$, then $q\alpha_{j,n}\in\mathbb Z$ for
$j=1,2,3$, while
\[
  |qt_{j,n}|\leq |q|t_{1,n}<\frac12.
\]
Thus $p_j(q)=q\alpha_{j,n}$ and
\begin{equation}\label{eq:cap-codim-one-special-return}
  \xi_j(q)=qt_{j,n}.
\end{equation}
Put
\begin{equation}\label{eq:cap-codim-one-rn}
  r_n:=\|q_n(t_{1,n},t_{2,n},t_{3,n})\|_\infty
  =q_nt_{1,n}
  \asymp\frac{M_n^3}{M_{n+1}}.
\end{equation}

\smallskip\noindent\emph{Hyperuniformity of the original interval.}
For
$0<\varepsilon<1/4$, define the sup-norm return set
\begin{equation}\label{eq:cap-codim-one-return-set}
  \mathcal R_\varepsilon
  :=\{q\in\mathbb Z\setminus\{0\}:
       \operatorname{dist}_\infty(q\alpha,\mathbb Z^3)\leq\varepsilon\}.
\end{equation}
Let $(\xi,\eta)\in\Lambda\setminus\{0\}$ satisfy
$\|\xi\|_2<\varepsilon$.  Then $q\neq0$: if $q=0$, the identities
$\xi_j=-p_j$ and $\varepsilon<1/4$ force $p=0$.  Hence
$q\in\mathcal R_\varepsilon$.  Moreover
\[
  |\eta|=|q-\xi_3|
  \geq |q|-|\xi_3|
  \geq |q|-\frac14
  \geq\frac{|q|}{2}.
\]
Thus \eqref{eq:cap-diffraction-formula} and
\eqref{eq:cap-codim-one-window-cancellation} give
\begin{equation}\label{eq:cap-codim-one-hu1-sum}
  \sigma_W(B_\varepsilon^*)
  \ll
  \sum_{q\in\mathcal R_\varepsilon}
    \frac{|\xi_3(q)|^2}{q^2}.
\end{equation}
Fix $n$ and assume
\begin{equation}\label{eq:cap-codim-one-epsilon-scale}
  \frac1{4M_{n+1}}\leq\varepsilon<\frac1{4M_n}.
\end{equation}

\smallskip\noindent\emph{Regime I: $\varepsilon\geq r_n/2$.}
Suppose first that $\varepsilon\geq r_n/2$.  For returns with $|q|\leq A_n$,
Equations~\eqref{eq:cap-codim-one-return-isolation} and
\eqref{eq:cap-codim-one-special-return} give $q=kq_n$ and
$\xi_j(q)=kq_nt_{j,n}$.  The condition $|\xi_1(q)|\leq\varepsilon$ implies
\[
  |k|\leq \frac{\varepsilon}{q_nt_{1,n}}.
\]
For every nonzero such $k$,
\[
  \frac{|\xi_3(kq_n)|^2}{(kq_n)^2}=t_{3,n}^2.
\]
Hence the contribution of $|q|\leq A_n$ to
\eqref{eq:cap-codim-one-hu1-sum} is
\begin{equation}\label{eq:cap-codim-one-small-return-contribution}
  \ll
  \left(1+\frac{\varepsilon}{q_nt_{1,n}}\right)t_{3,n}^2.
\end{equation}
Since $\varepsilon\geq r_n/2=q_nt_{1,n}/2$, division by $\varepsilon^3$
and the tail bounds give
\begin{equation}\label{eq:cap-codim-one-small-return-normalized}
  \frac{1}{\varepsilon^3}
  \left(1+\frac{\varepsilon}{q_nt_{1,n}}\right)t_{3,n}^2
  \ll
  \frac{t_{3,n}^2}{q_n^3t_{1,n}^3}
  \ll\frac1{M_n^9M_{n+1}^3}
  \longrightarrow0.
\end{equation}
If $q,q'\in\mathcal R_\varepsilon$, then
\[
  \operatorname{dist}_\infty((q-q')\alpha,\mathbb Z^3)
  \leq 2\varepsilon<\frac1{2M_n}.
\]
If in addition $0<|q-q'|\leq A_n$, then
\eqref{eq:cap-codim-one-return-isolation} gives
$q-q'\in q_n\mathbb Z$, and hence $|q-q'|\geq q_n$.  Thus distinct elements
of $\mathcal R_\varepsilon\cap\{|q|>A_n\}$ are $q_n$-separated.  Splitting the
positive and negative integers into intervals of length $q_n$ gives
\[
  \sum_{\substack{q\in\mathcal R_\varepsilon\\|q|>A_n}}q^{-2}
  \ll \sum_{r\geq0}(A_n+rq_n)^{-2}
  \ll\frac1{A_nq_n}.
\]
Using $|\xi_3(q)|<\varepsilon$, their contribution to
\eqref{eq:cap-codim-one-hu1-sum}, divided by $\varepsilon^3$, is
\begin{equation}\label{eq:cap-codim-one-large-return-normalized}
  \ll\frac1{\varepsilon A_nq_n}
  \ll\frac1{M_n^5},
\end{equation}
where the last estimate uses
$\varepsilon\geq r_n/2\asymp q_n/M_{n+1}$.

\smallskip\noindent\emph{Regime II: $\varepsilon<r_n/2$.}
Suppose next that $\varepsilon<r_n/2$.  Again $2\varepsilon<1/(2M_n)$.  If
$0<|q|\leq A_n$ is a $2\varepsilon$-return, then
\eqref{eq:cap-codim-one-return-isolation} gives $q=kq_n$, and
\eqref{eq:cap-codim-one-special-return} gives
\[
  \operatorname{dist}_\infty(q\alpha,\mathbb Z^3)
  =|k|r_n\geq r_n>2\varepsilon,
\]
a contradiction.  Hence every $\varepsilon$-return has $|q|>A_n$.
If two distinct $\varepsilon$-returns satisfied $|q-q'|\leq A_n$, their
difference would be a nonzero $2\varepsilon$-return in this forbidden range.
Thus $|q-q'|>A_n$, and
\[
  \sum_{q\in\mathcal R_\varepsilon}q^{-2}
  \ll\sum_{r\geq1}(rA_n)^{-2}
  \ll A_n^{-2}.
\]
and \eqref{eq:cap-codim-one-hu1-sum} gives
\begin{equation}\label{eq:cap-codim-one-second-regime}
  \frac{\sigma_W(B_\varepsilon^*)}{\varepsilon^3}
  \ll\frac1{\varepsilon A_n^2}
  \leq C\frac{M_n^2}{M_{n+1}}
  \longrightarrow0,
\end{equation}
because $\varepsilon\geq1/(4M_{n+1})$.  The two regimes cover every
sufficiently small $\varepsilon$.  Hence
\(
  \sigma_W(B_\varepsilon^*)=o(\varepsilon^3),
\)
and Theorem~\ref{thm:geometric-spectral} gives $\mu_W\in\HU_1$.

\smallskip\noindent\emph{Failure for a derived interval.}
The vector $\gamma_1$ from \eqref{eq:cap-codim-one-primal-vectors} has internal
component $\alpha_1\in(0,1)$, so
\begin{equation}\label{eq:cap-codim-one-derived-window}
  W':=W\cap(W-\alpha_1)
\end{equation}
is an interval of length $1-\alpha_1$.  For every $(\xi,\eta)\in\Lambda$,
\[
  (1-\alpha_1)\eta
  \equiv-\xi_1-(1-\alpha_1)\xi_3\pmod{\mathbb Z},
\]
and therefore
\begin{equation}\label{eq:cap-codim-one-derived-fourier}
  |\widehat{\chi}_{W'}(\eta)|
  =\frac{|\sin\pi(\xi_1+(1-\alpha_1)\xi_3)|}{\pi|\eta|}.
\end{equation}
At $q=q_n=M_n^3$, choose
\begin{equation}\label{eq:cap-codim-one-special-dual-p}
  p_{j,n}:=q_n\alpha_{j,n}\in\mathbb Z.
\end{equation}
Then
\[
  \xi_{j,n}=q_nt_{j,n}>0,
  \qquad
  \eta_n=q_n-\xi_{3,n},
\]
with $0<\xi_{3,n}<1$ for all sufficiently large $n$ and
$\|\xi_n\|_2\to0$.  The first coordinate dominates:
\begin{equation}\label{eq:cap-codim-one-first-coordinate-dominates}
  \|\xi_n\|_2\asymp\xi_{1,n}
  \asymp\frac{M_n^3}{M_{n+1}},
  \qquad
  \xi_{3,n}=o(\xi_{1,n}).
\end{equation}
Put
\[
  u_n:=\xi_{1,n}+(1-\alpha_1)\xi_{3,n}.
\]
By \eqref{eq:cap-codim-one-first-coordinate-dominates},
$u_n/\xi_{1,n}\to1$ and $u_n\to0$, while
$\eta_n/q_n\to1$.  Hence, for all large $n$,
$|\sin(\pi u_n)|\geq 2|u_n|$ and $|\eta_n|\leq2q_n$.  Equation
\eqref{eq:cap-codim-one-derived-fourier} therefore gives
\begin{equation}\label{eq:cap-codim-one-derived-atom-size}
  |\widehat{\chi}_{W'}(\eta_n)|^2
  \gg\frac{\xi_{1,n}^2}{q_n^2}.
\end{equation}
Taking $\varepsilon_n:=2\|\xi_n\|_2$ in
\eqref{eq:cap-diffraction-formula} gives
\begin{equation}\label{eq:cap-codim-one-derived-nonhu}
  \frac{\sigma_{W'}(B_{\varepsilon_n}^*)}{\varepsilon_n^3}
  \gg\frac1{q_n^2\xi_{1,n}}
  \asymp\frac{M_{n+1}}{M_n^9}
  \longrightarrow\infty.
\end{equation}
Hence $\mu_{W'}\notin\HU_1$ by
Theorem~\ref{thm:geometric-spectral}.

\smallskip\noindent\emph{From the derived window to failure of $\HU_2$.}
Finally, let
\begin{equation}\label{eq:cap-codim-one-D1}
  D_1:=-e_1+\alpha_1e_3
\end{equation}
be the physical component of $\gamma_1$.  The set
\(
  \{\gamma_V:\gamma\in\Gamma,\ \gamma_H\in W-W\}
\)
has a positive separation constant.  Under the identification
$\Sh_2(\mathbb R^3)\cong\mathbb R^3/\{D\sim-D\}$,
$D\mapsto[(-D/2,D/2)]$, its image is therefore discrete away from the
diagonal shape.  Choose
$\varphi\in C_c(\Sh_2(\mathbb R^3))$ equal to one at the shape of
$(-D_1/2,D_1/2)$ and supported so narrowly that no other nonzero displacement
class contributes; in particular, $\varphi(0)=0$.

For $\omega=\Gamma+(v,h)$, an accepted pair with displacement $D_1$ is
generated by $\gamma$ and $\gamma+\gamma_1$.  Both points occur exactly when
$\gamma_H+h\in W'$, and their barycentre is
$\gamma_V+v+D_1/2$.  The reversed ordering contributes a second ordered pair
with the same barycentre and the same unordered shape.  Therefore, exactly,
\begin{equation}\label{eq:cap-codim-one-pair-reduction}
  T_2[f,\varphi]
  =2S_{W'}(\tau_{D_1/2}f).
\end{equation}
This is the case $k=j=2$ and multiplicities $(1,1)$ of
\eqref{eq:cap-isolated-derived-window-statistic}.
If $\mu_W\in\HU_2$, applying the defining variance condition to this fixed
$\varphi$ and $f=\chi_{B_R}$ would imply
\[
  \Var\bigl(S_{W'}(\tau_{D_1/2}\chi_{B_R})\bigr)
  =o(\Vol_3(B_R)).
\]
Translation invariance of $\mu_{W'}$ removes the fixed translate of the ball,
so $\mu_{W'}\in\HU_1$, contradicting
\eqref{eq:cap-codim-one-derived-nonhu}.  Thus
$\mu_W\notin\HU_2$.
\end{proof}

\end{document}
